\documentclass[11pt]{article}
\usepackage[utf8]{inputenc}
\usepackage{amsmath,amssymb}
\usepackage[margin=1in]{geometry}
\usepackage{hyperref}
\emergencystretch=3em

\title{The density expansion of a face term}
\author{Claude, with Daniel Murfet}
\date{2026-09-07}

\begin{document}
\maketitle
\sloppy

\begin{abstract}
For a one-variable amplitude $\Psi(u,s)$ with $u$-Taylor data $f_m(s)$ to order $M$ and weights
$u^{h_1}v^{h_2}$, $p_2=(h_2+1)/k_2$, the density of unit~103 is expanded by the weighted
finite-part axis lemma (unit~100) with $\gamma=k_1p_2-h_1-1<M$:
\[
\rho(r,s)=\sum_{m<M}r^{p_2+\frac{m-\gamma}{k_1}-1}c_m(s)+r^{p_2-1}\,k_2^{-1}\mathrm{FP}_\gamma\Psi(\cdot,s)
+r^{p_2-1}\log\tfrac1r\,\frac{f_\gamma(s)}{k_1k_2}+O\!\bigl(r^{p_2+\frac{M-\gamma}{k_1}-1}H(s)\bigr),
\]
with the logarithm present exactly when $\gamma$ is an integer below $M$ (a collision of
exponents) and explicit $c_m$. The expansion is packaged as a \texttt{DensityExpansion} indexed by
$\mathrm{Fin}\,M\oplus\mathrm{Fin}\,2$ and transferred to the $N$-expansion of the face term. This
is the first concrete instance of Astra~\#4's cutoff-indexed $d=2$ density theorem: one face term,
all of its exponents, its resonant logarithm and its finite-part coefficient.
Zero \texttt{sorry}/\texttt{axiom}.
\end{abstract}

\section{The things formalised}
{\small
\begin{verbatim}
noncomputable def faceExp (p₂ γ : ℝ) (k₁ : ℕ) (m : ℕ) : ℝ := p₂ + ((m : ℝ) - γ) / k₁
noncomputable def faceCoeff (γ A : ℝ) (k₁ k₂ : ℕ) (f : ℕ → ℝ → ℝ) (m : ℕ) (s : ℝ) : ℝ :=
  if (m : ℝ) = γ then 1 / ((k₁ : ℝ) * k₂) * Real.log A * f m s
  else -(1 / (k₂ : ℝ)) * A ^ (-(((m : ℝ) - γ) / k₁)) / ((m : ℝ) - γ) * f m s
noncomputable def faceLogCoeff (γ : ℝ) (k₁ k₂ : ℕ) (f : ℕ → ℝ → ℝ) (M : ℕ) (s : ℝ) : ℝ :=
  ∑ m ∈ range M, if (m : ℝ) = γ then 1 / ((k₁ : ℝ) * k₂) * f m s else 0
noncomputable def faceFPCoeff (γ b : ℝ) (k₂ : ℕ) (Ψ : ℝ → ℝ → ℝ) (f : ℕ → ℝ → ℝ) (M : ℕ)
    (s : ℝ) : ℝ := 1 / (k₂ : ℝ) * axisFinitePart γ b (fun u => Ψ u s) (fun m => f m s) M
noncomputable def faceα / def faceJ / noncomputable def faceC
  -- indexed by Fin M ⊕ Fin 2
noncomputable def faceEnv (γ A : ℝ) (k₁ k₂ M : ℕ) (H : ℝ → ℝ) (s : ℝ) : ℝ :=
  1 / (k₂ : ℝ) * (A ^ (-(((M : ℝ) - γ) / k₁)) / ((M : ℝ) - γ)) * H s
theorem face_monomial_eq ... :
    -(1 / (k₂ : ℝ)) * r ^ (p₂ - 1) * (f m s * axisPrim γ ((r / b ^ k₂) ^ ((k₁ : ℝ)⁻¹)) m)
      = r ^ (faceExp p₂ γ k₁ m - 1) * faceCoeff γ (b ^ k₂) k₁ k₂ f m s
        + r ^ (p₂ - 1) * Real.log (1 / r)
          * (if (m : ℝ) = γ then 1 / ((k₁ : ℝ) * k₂) * f m s else 0)
theorem face_terms_eq ...  -- local expression = ∑ᵢ r^{αᵢ-1} log(1/r)^{jᵢ} cᵢ(s)
theorem genDensity_measurable / faceFPCoeff_measurable / faceC_measurable
theorem faceDensity_expansion (p₂ b : ℝ) (h₁ k₁ k₂ M : ℕ) (hb : 0 < b) (hk₁ : 0 < k₁)
    (Ψ : ℝ → ℝ → ℝ) (hΨc : Continuous (Function.uncurry Ψ)) (f : ℕ → ℝ → ℝ)
    (hf : ∀ m, Measurable (f m)) (H : ℝ → ℝ) (hH : Measurable H) (hH0 : ∀ s, 0 ≤ H s)
    (hγ : (k₁ : ℝ) * p₂ - h₁ - 1 < M)
    (hrem : ∀ s, 0 < s → ∀ u ∈ Ioc (0 : ℝ) b,
      |Ψ u s - ∑ m ∈ range M, f m s * u ^ m| ≤ H s * u ^ M) :
    DensityExpansion (genDensity p₂ b h₁ k₁ k₂ Ψ) (faceα p₂ γ k₁ M) (faceJ M)
      (faceC γ b k₁ k₂ Ψ f M) (p₂ + ((M : ℝ) - γ) / k₁) (b ^ (k₁ + k₂)) 1
      (faceEnv γ (b ^ k₂) k₁ k₂ M H)   -- γ = k₁p₂-h₁-1
theorem face_transfer_bound ... (hp₂ : ((h₂ : ℝ) + 1) / k₂ = p₂) ... (N : ℝ) (hN : 1 ≤ N) :
    |twoDAmp β b N h₁ h₂ k₁ k₂ (fun u _ s => Ψ u s)
        - ∑ i, transferTerm β (faceα ..) (faceJ ..) (faceC ..) N|
      ≤ N ^ (-(p₂ + (M - γ)/k₁)) * (1 + Real.log N) ^ 1
        * (envMoment .. + ∑ i, R ^ (αᵢ - a) * tailMoment ..)
theorem face_transfer_sum_eq ... :
    ∑ i, transferTerm β (faceα p₂ γ k₁ M i) (faceJ M i) (faceC γ b k₁ k₂ Ψ f M i) N
      = (∑ m : Fin M, N ^ (-(faceExp p₂ γ k₁ m))
          * logMoment β (faceExp p₂ γ k₁ m) 0 (faceCoeff γ (b ^ k₂) k₁ k₂ f m))
        + N ^ (-p₂) * logMoment β p₂ 0 (faceFPCoeff γ b k₂ Ψ f M)
        + N ^ (-p₂) * (Real.log N * logMoment β p₂ 0 (faceLogCoeff γ k₁ k₂ f M)
          - logMoment β p₂ 1 (faceLogCoeff γ k₁ k₂ f M))
\end{verbatim}
}

\section{Mathematics}

With $\varepsilon(r)=(r/b^{k_2})^{1/k_1}$ and $A=b^{k_2}$, unit~100 gives
$J(\varepsilon)=\mathrm{FP}_\gamma-\sum_{m<M}f_m\,\mathrm{axisPrim}(\gamma,\varepsilon,m)+E$ with
$|E|\le H\varepsilon^{M-\gamma}/(M-\gamma)$. The monomial terms are converted to powers of $r$ by
\texttt{axisPrim\_rpow}: off resonance $(r/A)^{(m-\gamma)/k_1}/(m-\gamma)=r^{(m-\gamma)/k_1}A^{-(m-\gamma)/k_1}/(m-\gamma)$,
so the term has exponent $p_2+(m-\gamma)/k_1$; at the resonance $m=\gamma$,
$(\log r-\log A)/k_1$ splits into a $\log(1/r)$ term with coefficient $(k_1k_2)^{-1}f_\gamma$ and a
constant $(k_1k_2)^{-1}\log A\,f_\gamma$ at exponent $p_2$ (\texttt{face\_monomial\_eq}). The
remainder $k_2^{-1}r^{p_2-1}\varepsilon^{M-\gamma}H/(M-\gamma)$ has exponent
$a=p_2+(M-\gamma)/k_1$, strictly above every term since $m<M$; the uniform log degree $1$ absorbs
the resonant term. Measurability of the finite-part coefficient in $s$ is a parametric-integral
argument (\texttt{StronglyMeasurable.integral\_prod\_right'}), and of the density in $(r,s)$ the
indicator trick of unit~87.

\section{Paper-facing meaning}

A face term $u^ia_i(v,s)$ of the rectangular decomposition (after swapping coordinates) is exactly
this situation with $\gamma=k_2i/k_1$, so its contribution to the $d=2$ block is
\begin{gather*}
\sum_{m<M}N^{-(p+m/k_2)}\,\mathcal M_{p+m/k_2,0}(c_m)
+N^{-(p+i/k_1)}\Bigl[\mathcal M_{p+i/k_1,0}(k_2^{-1}\mathrm{FP})\\
\qquad+\log N\,\mathcal M_{p+i/k_1,0}(c_{\log})-\mathcal M_{p+i/k_1,1}(c_{\log})\Bigr]
+O\bigl(N^{-(p+M/k_2)}(1+\log N)\bigr),
\end{gather*}
with $c_{\log}\ne0$ only when $i/k_1=m/k_2$ for some $m<M$: the logarithms of the higher-order
$d=2$ expansion occur precisely at collisions of the two exponent lattices $p+\mathbb N/k_1$ and
$p+\mathbb N/k_2$, as predicted in Astra~\#3(b)/\#4(b). What remains for the full $\tau$-cutoff
theorem is the rectangular decomposition $\Phi=\sum_iu^ia_i+\sum_jv^jb_j-\sum_{ij}c_{ij}u^iv^j+R$
with the mixed remainder estimate, and the bookkeeping of the corner monomials.

\section{Lean notes}

\begin{itemize}
\item A file's import closure need not contain everything below it in \texttt{Laplace.lean}:
\texttt{axisFinitePart} was unknown until \texttt{AxisSeries} (hence \texttt{FinitePartAxis}) was
imported explicitly.
\item \texttt{Fintype.sum\_sum\_type} + \texttt{Fin.sum\_univ\_two} + \texttt{simp only
[Sum.elim\_inl, Sum.elim\_inr, Matrix.cons\_val\_zero, Matrix.cons\_val\_one]} unfold a
\texttt{Fin M ⊕ Fin 2}-indexed sum; \texttt{Fin.sum\_univ\_eq\_sum\_range} aligns the
\texttt{Fin M} sum with a \texttt{range M} sum; the final step is
\texttt{linear\_combination (-1 : ℝ) * hS}.
\item \texttt{simp only} may already distribute a scalar over a sum, so a later
\texttt{rw [Finset.mul\_sum]} finds nothing; prove the combined identity as a \texttt{have} instead.
\item A \texttt{gcongr} side goal \texttt{0 ≤ faceEnv …} is reported as an unsolved
\texttt{case ha} at the enclosing \texttt{have}; opaque definitions need the monotonicity lemma
spelled out (\texttt{mul\_le\_mul\_of\_nonneg\_right}).
\item \texttt{split\_ifs} on a condition independent of the variable (\texttt{(m : ℝ) = γ}) is the
cheap way to get measurability of an \texttt{ite}-defined coefficient.
\end{itemize}

\end{document}
